% Generated draft: requires mathematical and source review.
\documentclass[11pt]{article}
\usepackage[margin=25mm]{geometry}
\usepackage{fontspec}
\setmainfont{Latin Modern Roman}
\usepackage{amsmath,amssymb,mathtools,unicode-math}
\setmathfont{Latin Modern Math}
\usepackage{enumitem,needspace,xurl,hyperref}
\setlength{\parindent}{0pt}
\setlength{\parskip}{5pt}
\setlength{\emergencystretch}{3em}
\newcommand{\N}{\mathbb{N}}
\newcommand{\Z}{\mathbb{Z}}
\newcommand{\Q}{\mathbb{Q}}
\newcommand{\R}{\mathbb{R}}
\newcommand{\C}{\mathbb{C}}
\newcommand{\F}{\mathbb{F}}
\newcommand{\A}{\mathbb{A}}
\newcommand{\PP}{\mathbb{P}}
\newcommand{\E}{\mathbb{E}}
\newcommand{\eps}{\varepsilon}
\newcommand{\dd}{\,\mathrm{d}}
\begin{document}
\section*{UnsolvedMath ID: 30001988}
\textbf{Source problem code:} catalogue entry
\subsection{Definitions and mathematical statement}
A computable differential field is a field with a computable presentation in which the field operations and the distinguished derivation are computable. In the ordinary differential-field convention, a derivation \(\delta\) is an additive map satisfying \(\delta(ab)=a\delta(b)+b\delta(a)\). The supplied record does not explicitly specify the characteristic or the number of derivations; the notation below unpacks the ordinary, single-derivation interpretation without imposing characteristic zero.

Let \((K,\delta)\) be a computable differential field contained in a larger computable differential field \((L,\delta)\), with the derivation on \(L\) restricting to that on \(K\). An element \(z\in L\) is differentially transcendental over \(K\) if, for every integer \(n\geq 0\) and every nonzero polynomial \(P\in K[X_0,\ldots,X_n]\),
\[
P(z,\delta z,\ldots,\delta^n z)\neq 0.
\]
Write
\[
K\langle z\rangle=K(z,\delta z,\delta^2z,\ldots)
\]
for the smallest differential subfield of \(L\) containing \(K\) and \(z\).

The supplied question asks whether, for every such \(K\), \(L\), and \(z\),
\[
T_{K\langle z\rangle}\leq_T T_K.
\]
For subsets \(A,B\subseteq\N\), the relation \(A\leq_T B\) means that an oracle Turing machine with oracle \(B\) computes the characteristic function of \(A\).

\textbf{Unresolved qualification.} The supplied record does not define \(T_F\), its effective coding, or its dependence on a presentation of the differential field \(F\). Thus the displayed question preserves the source notation but is not yet a fully specified computability statement. The title mentions splitting sets, but this alone does not identify \(T_F\). The record also does not state whether the inclusion \(K\subseteq L\) must be computable or how \(K\langle z\rangle\) is to be presented effectively. No uniform reduction across all instances is explicitly requested.
\subsection{Short English statement}
The question asks whether adjoining a differentially transcendental element to a computable differential field preserves computability of the associated set \(T_F\) relative to its value for the original field. The meaning and coding of \(T_F\) are missing, so a precise formulation requires further source clarification. The supplied curation labels the problem open and reports that a relevant proof remains to be checked; these are supplied assessments, not independently verified conclusions.
\subsection{Sources}
\begin{itemize}
\item Supplied local UnsolvedMath registry record, permanent ID \(30001988\), exact source problem code \texttt{catalogue entry}, title \emph{Splitting Sets Under Differential-Transcendental Extensions}; collection: Oberwolfach Reports Open Problems. The record gives the proposed year as 2012 and supplies no individual proposer.
\item \emph{Computability Theory} (2012), Oberwolfach Reports, pp.\ 397--455, DOI \texttt{10.4171/owr/2012/08}. Supplied source URL and literature reference labeled ``original Oberwolfach report'': \url{https://doi.org/10.4171/owr/2012/08}.
\item Literature assessment dated \texttt{2026-08-21}: the supplied curation labels the problem open and states that the report presents the statement as a conjecture whose relevant proof remains to be checked. This is supplied dated curation, not independently verified literature assessment; the record itself requests independent verification of time-sensitive status and cited claims.
\item The registry labels the statement \texttt{corrected\_verified} and reports comparison with the official Oberwolfach report, with notation, context, and question form normalized without changing the source-level claim. This is supplied verification metadata; no independent source inspection or mathematical verification is asserted here.
\end{itemize}
\end{document}
